\section{Discussion}
\label{sec:discussion}
\textbf{Core challenges in LLM agent security defense.} 
A central motivation of this work is that existing agent security frameworks overwhelmingly rely on static rules, fixed domain knowledge, or predefined tools, leading to poor performance when confronted with dynamically open environments. We propose DEFEND, the first framework to integrate agent self-evolution capabilities with security defense mechanisms, treating ``defense'' itself as a generable, continuously optimizable artifact with lifelong learning capabilities. Experimental results (Table~\ref{tab:main_results}) validate the effectiveness of this paradigm shift: DEFEND maintains stable and superior defensive performance across benchmarks and bacebone LLMs, while baseline methods exhibit drastic performance fluctuations of up to 20\%.

\textbf{Effectiveness of tool generation and self-evolution mechanisms.}
A core design decision in DEFEND is to translate abstract security policies into executable Python guardrail tool code, rather than relying on natural language discrimination logic. This decision has a deep epistemological basis. Natural language discrimination mechanisms inherently suffer from ambiguity, manipulability, and interpretive inconsistency, whereas code-based guardrail tools provide a more rigorous and verifiable protection logic. ``Codifying'' security policies not only makes the protection logic more precise, but also more portable: the same tool can be reused across different backbone models and application scenarios without redesigning the defense logic for each case.

In terms of the self-evolution mechanism, DEFEND achieves knowledge accumulation and continuous optimization at the tool level through FusionAgent. Results (Figure~\ref{fig:reuse_result}) show that DEFEND can dynamically determine when guardrail tools are needed, achieving tool reuse rates exceeding 53\% across multiple domains. The significance of this mechanism extends beyond reducing redundant work. Its deeper value lies in gradually building a tool knowledge library spanning multiple domains and risk types as interaction experience accumulates. Each successful defense experience is preserved as a reusable tool, and each tool fusion represents an expansion and deepening of existing protection capabilities. This accumulation mechanism closely mirrors the human process of summarizing experience and refining response strategies through practice, embodying a genuinely cognitive evolutionary logic.

\textbf{The VGT phenomenon and the necessity of auditing.} 
A notable finding of this work is a previously underappreciated phenomenon: even when safety constraints are explicitly incorporated into tool generation prompts, the guardrail tools initially produced by agents may still harbor security vulnerabilities, a phenomenon we define as Vulnerabilities in Guardrail Tools (VGT). When agents are tasked with generating guardrail tools, their attention tends to concentrate heavily on whether the tool can effectively detect target risks, while paying insufficient attention to the code security of the tool itself. This is essentially the same as the "functionally correct but vulnerable" phenomenon in software engineering. Therefore, newly generated knowledge requires secondary validation, and security should not be understood as a dynamic process requiring continuous verification throughout the lifecycle.

It is precisely based on this recognition that DEFEND introduces an independent auditing mechanism. AuditorAgent does not participate in the tool generation and fusion process, and is able to conduct security reviews of newly generated or fused tools with an objective and conservative stance. By decoupling tool verification from final decision-making, AuditorAgent serves as a safety anchor across the entire defense chain: even if upstream modules introduce potential defects during tool generation or optimization, the auditing mechanism can still block unsafe tools from entering the tool library based on execution evidence and conservative auditing principles, thereby effectively suppressing the misevolution risk caused by VGT phenomenon.

\textbf{Theoretical implications of the ablation study.}
The ablation results not only validate the necessity of each core component, but also reveal the internal consistency of DEFEND's overall design logic from another angle. The absence of each component corresponds to a distinct dimension of degradation: removing AnalysisAgent leads to a significant decline in cross-domain generalization, indicating that dynamic risk perception is the core driver for transcending domain boundaries; the long-term effect of removing FusionAgent manifests as reduced system efficiency and sustainability; removing AuditorAgent results in a notable drop in recall, revealing that without independent audit constraints, excessive reliance on tool execution results leads to misjudgments of benign samples.

These three degradation patterns are not independent of one another, but are deeply interrelated. The absence of dynamic risk analysis means that generated tools cannot precisely target emerging threats, thereby undermining the quality of tool fusion; the absence of the tool fusion mechanism traps knowledge accumulation in an inefficient loop; and the absence of the auditing mechanism leaves the entire self-evolution process without credible quality assurance. The synergy of these three components jointly constitutes a complete closed loop spanning risk perception, tool evolution, and trusted verification. Security is not the simple superposition of point-by-point protections, but an emergent property of multiple mechanisms mutually constraining and co-evolving with one another.

\textbf{Limitations and future research directions.}
Dynamic tool generation, sandboxed execution, and auditing mechanisms inevitably introduce additional time overhead compared to pure natural language discrimination or static rule filtering. In application scenarios with extremely high real-time demands (such as real-time gaming, or high-frequency trading), this overhead may pose a practical barrier to deployment. Although FusionAgent's tool reuse mechanism can partially amortize the cost of tool generation, how to effectively control overhead while maintaining defense quality as the system operates in increasingly complex environments remains an open question requiring further exploration. Looking ahead, investigating how to better apply DEFEND to tasks such as robot control and real-time gaming, and exploring how to achieve distributed, collaborative self-security evolution in multi-agent cooperative or competitive settings, represents a more cutting-edge and promising direction.

\textbf{Conclusion.}
In this paper, we propose DEFEND, the first dynamic defense framework that integrates self-evolution capabilities with agent security defense mechanisms, addressing the limitations of existing frameworks in terms of poor generalization and lack of adaptive capability. DEFEND transforms abstract security policies into generable, executable, and continuously optimizable guardrail tool code, making the defense mechanism itself an evolving artifact with lifelong learning capabilities. We further discover and define the phenomenon of Vulnerabilities in Guardrail Tools (VGT), revealing the intrinsic risks that self-evolving methods face during tool generation, and address this through an independent auditing mechanism. DEFEND demonstrates stable and significant defensive advantages across multiple benchmark datasets and models, validating the feasibility of building a general-purpose agent defense system driven by self-evolution at its core.

% This study introduces DEFEND, the first dynamic defense framework that integrates agent self-evolution capabilities with agent security protection. Through its core four agent components, the framework translates abstract security requirements into executable, verifiable, and optimizable code tools, enabling continuous evolution throughout the system lifecycle. Experimental results demonstrate that DEFEND exhibits outstanding defensive performance, strong cross-domain and cross-model generalization capabilities, and high trustworthiness ensured by end-point auditing mechanisms across multiple benchmarks and base models. This section will delve into the significance, limitations, and future directions of this work.

% Traditional agent defense mechanisms predominantly rely on static rules, predefined tools, or domain-specific knowledge, which severely limits their generalization capability when confronted with dynamically evolving open environments. The core contribution of DEFEND framework lies in its novel conceptualization: it treats ``defense'' itself as a dynamic, generable, continuously optimizable, and lifelong-learnable artifact. By emulating the human process of self-evolution, DEFEND can proactively abstract security patterns from interaction with the environment, rather than passively depending on predefined rules. Experimental results (shown in Figure~\ref{fig:main_result}) clearly validate the superiority of this paradigm: DEFEND maintains stable and superior performance across diverse benchmarks with distinct properties, significantly outperforming baseline methods that exhibit highly volatile performance in cross-benchmark tests. This marks a shift in agent defense from a static patch requiring manual priors and domain adaptation to a general framework endowed with adaptability and dynamism. Furthermore, the auditing mechanism prevents the self-evolution process from deviating. Through its AuditorAgent component, DEFEND subjects every newly generated or fused tool to a re-questioning of its logical integrity and security. This ensures that tools stored in the tool library for reuse possess high reliability, providing a security guarantee for building truly trustworthy and secure dynamic agent systems.

% Dynamic generation, simulated execution of Python security tools, and multi-round auditing inevitably introduce additional computational and time overhead compared to pure natural language discrimination or static rule filtering. Although the FusionAgent effectively mitigates knowledge library bloat through tool reuse, the current framework's latency may require further optimization for scenarios with extremely high real-time demands. Therefore, for future directions, researching how to better apply DEFEND to tasks such as robot control and real-time gaming, and exploring how to achieve distributed, collaborative self-security evolution, presents a more cutting-edge and promising path.

% In summary, as a pioneering exploration towards a unified framework for agent defense, DEFEND, by making security mechanisms tool-based, dynamic, and evolvable, points to a new direction for building more general, adaptive, and trustworthy agent protection systems. It not only provides an effective solution to current security challenges faced by agents but also offers a feasible approach for constructing future intelligent systems that can safely and autonomously learn and adapt to the complex real world.